% ===== §1 Introduction =====
\section{Introduction: From the Museum to the Kitchen}
\label{p22:sec:introduction}

\epigraph{\zh{担水砍柴，无非妙道。}}{\textit{Chan saying}}

\noindent
Aesthetic education has, for most of its history, taken place before the masterwork. One is brought before the painting, the symphony, the poem, and one's taste is refined by the encounter, so that the educated person is the one who perceives finely what is finely made. This paper proposes to move aesthetic education from the museum to the kitchen, from the masterwork to the morning, and the proposal is not a populist gesture, a democratisation of access to the fine, but a claim about where the work of aesthetic education most fundamentally lies. The claim is that the deepest aesthetic education is not the refinement of taste before objects but the cultivation of the capacity to build generative fields between people, and that this capacity is exercised and acquired not in the exceptional encounter with the masterwork but in the most ordinary and repetitive acts of a shared life, in the placing of a cup and the making of a meal and the meeting of a morning's small contingency. Why ethics should begin from washing the dishes, why the kitchen rather than the museum should be the privileged site of an education in beauty, is the paradox the paper exists to resolve.

\medskip\noindent

The paradox is sharp because the museum and the kitchen seem to lie at opposite ends of any scale of the aesthetic, the one the home of beauty and the other the home of drudgery, and to locate the deepest aesthetic education in the kitchen seems either a sentimentality or a confusion. The paper resolves the paradox by a chain of argument that the introduction now previews, so that the reader knows the shape of the road before walking it. Relational production, the generation of the relational subject established in the immediately preceding work, requires a field, a conditional environment in which the generative coupling can occur and accumulate. The construction of such a field, the paper will argue, is not an engineering matter, because the appropriateness it must achieve is state-dependent and cannot be specified, and so it is an aesthetic matter, a discernment of the fitting where no rule decides. The aesthetic, pressed to its ground, discloses the ethical, the two coupled in a single perception of the field's generative direction. And so the cultivation of the capacity to build generative fields, which is a cultivation of this coupled perception, is what aesthetic education most fundamentally is. The chain runs from relational production to the field, from the field to aesthetics, from aesthetics to ethics, and from there to aesthetic education, and its terminus is the resolution of the paradox: the kitchen is the privileged site because the field is built there, in the ordinary acts of the shared life, and the building of it is the aesthetic education the museum only ever partially displayed.

\medskip\noindent

The paper makes several contributions that the introduction names so that they may be watched for. It establishes a concept, relational aesthetics, on which beauty is a coupling event between subjects rather than a property of objects or a state of subjects, and it argues that this concept is required by the relational ontology the series has carried throughout. It proposes a model, the four movements of perception, reception, sharing, and reproduction, and it shows the model to be not a procedure laid over experience but the smallest form a good cycle takes when enacted between two people, a holonomy cycle that returns to its root and not its origin. It reconstructs a method, the dialectic of principal and secondary contradiction, into a generative form that directs attention rather than concentrating force. It gathers, at four points where the model's movements carry the densest theory, the contending traditions of East and West, and it shows them resolved not by a synthesis that absorbs them but by an arrangement that keeps their differences while finding their common ground. And it faces, rather than evading, the political charge that an aesthetics of the everyday is a doctrine for the leisured, arguing that the cultivation it describes requires no surplus of time or means and so is available precisely to those whose days hold no surplus, the dispossessed for whom the difference between the generative and the deadening repetition is most consequential of all.

\medskip\noindent

This paper occupies a particular place in the series it belongs to, and naming the place orients the reader who comes to it from the earlier work. The series has proceeded largely by studying particular fields, the conditional environments in which the good cycle is secured, each studied for the way its conditions hold off the slide toward the bad: the field of poetry and its refusal of closure, the field of the vow and its self-legislation, the field of yielding water, the field of travel and its spaciousness. Each such study described a generative field and the way its particular conditions secured generativity, and each presupposed, without naming, a capacity to build and sustain such fields. This paper names that capacity and makes it the subject, and in doing so it stands to the prior studies as an account of the builder to accounts of the things built. Where the prior papers asked what a generative field is and how its conditions work, this paper asks how a person comes to be able to build generative fields at all, and answers that the coming-to-be-able is an aesthetic education, the cultivation of the perception and practice and maintenance by which fields are built. The paper is therefore late in the series by design, presupposing the studies of particular fields and gathering from them the question they share, the question of the capacity, which only a study devoted to the capacity itself could answer.

The relation to the immediately preceding work is closer still and should be marked. That work established relational production, the process by which the relational subject is generated, and located it in the field of travel under the figure of spaciousness. This paper takes up relational production where that work left it and asks what relational production requires in general, across all fields and not only the spacious field of travel, and finds that it requires a field, and that the building of fields is the aesthetic education the paper describes. The move from the preceding work to this one is the move from a particular field, travel, in which relational production was first formally established, to the general account of the field and its construction, of which the field of travel was one spacious instance. This paper thus generalises the preceding one, carrying relational production out of the spacious field of travel into the dense and repetitive fields of the everyday, and showing that the capacity to build generative fields, displayed in travel where spaciousness made it legible, is exercised and acquired in the unspacious everyday where most of relational production must in fact occur.

\medskip\noindent

The wager of the paper is stated plainly at the outset, because everything depends on it. The wager is that the good cycle can be built in the everyday, in the most repetitive and unspacious life, and not only in the heightened fields of the exceptional. An aesthetic education that worked only on the masterwork, or only in the spacious occasions of travel and ceremony, would be an education for the hours a life mostly does not have, leaving the bulk of a life, its mornings and repetitions and narrow rooms, uncultivated. The paper wagers that the everyday is not the barren ground from which beauty must be imported but the very site where the generative field is most needed and, properly understood, most buildable, and that the cultivation of the capacity to build it there is the aesthetic education that matters most, because a life is mostly everyday, and to make the everyday generative is to make most of a life so. The Chan saying at the head of this introduction states the wager in four words: carrying water and chopping firewood, the most ordinary and repetitive labour, is nothing other than the wondrous way. The paper is the working-out of that ancient claim in the terms of relational production, the field, and the holonomy that a good cycle accumulates, and its argument is, in the end, that the wondrous way is indeed in the carrying of the water, if the carrying is perceived and received and shared and renewed with the trained and unattached attention that the aesthetic education of the everyday exists to cultivate.
